\section{Python 源程序}
\renewcommand{\thesubsection}{A\thinskip.\thinskip\arabic{subsection}}
\subsection{问题一}
\vspace{-2ex}

\begin{Python}{Q1}
# ================= 数据分段 =================
def segment_signal(x, window_size=1024, overlap=0.5):
    """将信号按窗口切分"""
    step = int(window_size * (1 - overlap))
    segments = []
    for i in range(0, len(x) - window_size + 1, step):
        segments.append(x[i:i + window_size])
    return segments

# ================= 时域特征 =================
def time_features(x):
    """提取时域特征"""
    x = np.array(x)
    mean_val = np.mean(x)                         # 均值
    std_val = np.std(x)                           # 标准差
    var_val = np.var(x)                           # 方差
    rms = np.sqrt(np.mean(x**2))                  # 均方根
    peak = np.max(np.abs(x))                      # 峰值
    crest_factor = peak / rms if rms != 0 else 0  # 峭度因子
    impulse_factor = peak / np.mean(np.abs(x))    # 脉冲因子
    margin_factor = peak / (np.mean(np.sqrt(np.abs(x)))**2)  # 裕度因子
    kurt_val = kurtosis(x)                        # 峭度
    skew_val = skew(x)                            # 偏度
    return [mean_val, std_val, var_val, rms, peak, crest_factor,
            impulse_factor, margin_factor, kurt_val, skew_val]

# ================= 频域特征 =================
def freq_features(x, fs=12000):
    """提取频域特征，包括功率谱和频谱峭度"""
    nperseg = min(1024, len(x))  # 防止窗口过大
    f, Pxx = welch(x, fs=fs, nperseg=nperseg)
    Pxx = Pxx / np.sum(Pxx)  # 归一化
    centroid = np.sum(f * Pxx)                             # 频率质心
    peak_freq = f[np.argmax(Pxx)]                          # 峰值频率
    bandwidth = np.sqrt(np.sum(((f - centroid)**2) * Pxx)) # 频带宽度
    spectral_kurt = kurtosis(Pxx) if len(Pxx) > 3 else 0  # 避免太短的向量计算失败  # 频谱峭度
    return [centroid, peak_freq, bandwidth, spectral_kurt]

# ================= 包络谱特征 =================
def envelope_features(x, fs=12000):
    """Hilbert 包络谱特征"""
    analytic_signal = hilbert(x)
    envelope = np.abs(analytic_signal)
    f, Pxx = welch(envelope, fs=fs, nperseg=1024)
    Pxx = Pxx / np.sum(Pxx)
    peak_freq = f[np.argmax(Pxx)]
    energy = np.sum(Pxx)
    spectral_kurt = kurtosis(Pxx)
    return [energy, peak_freq, spectral_kurt]

# ================= 小波包特征 =================
def wpt_features(x, wavelet='db4', level=3):
    """小波包分解特征"""
    wp = pywt.WaveletPacket(data=x, wavelet=wavelet, maxlevel=level)
    nodes = wp.get_level(level, order='freq')
    energies = [np.sum(np.square(n.data)) for n in nodes]
    total_energy = np.sum(energies)
    energy_ratio = [e / total_energy for e in energies]
    entropy = -np.sum([e*np.log(e+1e-12) for e in energy_ratio])
    return energies + energy_ratio + [entropy]

# ================= STFT 特征 =================
def stft_features(x, fs=12000, nperseg=256):
    """短时傅里叶变换特征"""
    f, t, Zxx = stft(x, fs=fs, nperseg=nperseg)
    magnitude = np.abs(Zxx)
    mean_mag = np.mean(magnitude)
    std_mag = np.std(magnitude)
    max_mag = np.max(magnitude)
    return [mean_mag, std_mag, max_mag]

# ================= Hilbert–Huang 特征 =================
def hht_features(x):
    """Hilbert-Huang变换特征（简化版）"""
    analytic_signal = hilbert(x)
    amplitude_envelope = np.abs(analytic_signal)
    instantaneous_phase = np.unwrap(np.angle(analytic_signal))
    instantaneous_freq = np.diff(instantaneous_phase) / (2.0*np.pi)
    return [np.mean(amplitude_envelope), np.std(amplitude_envelope),
            np.mean(instantaneous_freq), np.std(instantaneous_freq)]

# ================= 轴承特征频率 =================
def bearing_frequencies(rpm, n, d, D):
    """计算轴承特征频率"""
    fr = rpm / 60.0
    BPFO = (n / 2) * fr * (1 - d / D)
    BPFI = (n / 2) * fr * (1 + d / D)
    BSF = (D / d) * fr * (1 - (d / D)**2)
    FTF = (1 / 2) * fr * (1 - d / D)
    return fr, BPFO, BPFI, BSF, FTF

# ================= 根据路径判断轴承型号 =================
def get_bearing_params_from_path(file_path):
    if "FE" in file_path:
        n, d, D = 9, 0.2656, 1.122
    elif "DE" in file_path:
        n, d, D = 9, 0.3126, 1.537
    elif "BA" in file_path:
        n, d, D = 9, 0.2656, 1.122
    elif "Normal" in file_path:
        n, d, D = 9, 0.3126, 1.537
    else:
        raise ValueError(f"Cannot determine bearing type from path: {file_path}")
    return n, d, D

# ================= 获取 RPM =================
def get_rpm_from_filename(file_path):
    match = re.search(r'\((\d+)rpm\)', file_path)
    if match:
        return float(match.group(1))
    else:
        return 1800.0

# ================= 读取信号 =================
def get_signals_from_file(file_path):
    data = scio.loadmat(file_path)
    signals = {}
    for ch in ["FE", "DE", "BA"]:
        keys = [k for k in data.keys() if k.endswith(f"_{ch}_time") and not k.startswith("__")]
        if keys:
            signals[ch] = data[keys[0]].squeeze()
    if not signals:
        raise ValueError(f"No vibration signals found in {file_path}")
    rpm_keys = [k for k in data.keys() if "RPM" in k and not k.startswith("__")]
    if rpm_keys:
        rpm_data = data[rpm_keys[0]]
        rpm = float(rpm_data.squeeze())
    else:
        rpm = get_rpm_from_filename(file_path)
    return signals, rpm

# ================= 标签解析 =================
def parse_label(file_path):
    if "OR" in file_path:
        return "OR"
    elif "IR" in file_path:
        return "IR"
    elif "B" in file_path:
        return "B"
    elif "N" in file_path:
        return "N"
    else:
        return "Unknown"

def parse_label_simple(file_path):
    """仅提取故障类型 OR / IR / B / N，用于窗口选择"""
    if "/B/" in file_path or "_B" in file_path:
        return "B"
    elif "/IR/" in file_path or "_IR" in file_path:
        return "IR"
    elif "/OR/" in file_path or "_OR" in file_path:
        return "OR"
    else:
        return "N"

def parse_label_full(file_path, label):
    """生成复合标签：通道_故障类型_故障尺寸_载荷"""
    # ===== 通道 =====
    if "FE" in file_path:
        channel = "FE"
    elif "DE" in file_path:
        channel = "DE"
    elif "BA" in file_path:
        channel = "BA"
    else:
        channel = "NO"
    # ===== 故障类型 =====
    fault_type = parse_label_simple(file_path)
    # ===== 故障尺寸 =====
    match_size = re.search(r'/(\d{4})/', file_path)
    fault_size = match_size.group(1) if match_size else "Unknown"

    # ===== 载荷 =====
    match_load = re.search(r'_([0-3])\.mat$', file_path)
    load = match_load.group(1) if match_load else "Unknown"

    if load == "Unknown":
        m = re.search(r'[ _]([0-3])(?:_\(\d+rpm\))?\.mat$', file_path)
        if m:
            load = m.group(1)

    print(fault_size)
    return f"{channel}_{fault_type}_{fault_size}_{load}_{label}"


# ================= 提取单文件特征 =================
def extract_features_from_file(file_path, fs=12000, overlap=0.5):
    signals, rpm = get_signals_from_file(file_path)
    label = parse_label(file_path)
    all_features = []
    labels = []
    n, d, D = get_bearing_params_from_path(file_path)

    for ch, sig in signals.items():
        rpm_mean = np.mean(rpm)
        fr, BPFO, BPFI, BSF, FTF = bearing_frequencies(rpm_mean, n, d, D)

        print(BPFO)

        if label == "OR":
            window_size = int(fs / BPFO)
        elif label == "IR":
            window_size = int(fs / BPFI)
        elif label == "B":
            window_size = int(fs / BSF)
        else:
            window_size = int(fs / 120)


        print(window_size)

        segments = segment_signal(sig, window_size, overlap)
        full_label = parse_label_full(file_path, ch)   # ← 移到这里
        print(full_label)
        for seg in segments:
            feat_t = time_features(seg)
            feat_f = freq_features(seg, fs)
            feat_env = envelope_features(seg, fs)
            feat_wpt = wpt_features(seg)
            feat_stft = stft_features(seg, fs)
            feat_hht = hht_features(seg)
            feat_bearing = [fr, BPFO, BPFI, BSF, FTF]
            feature_vec = [ch] + feat_t + feat_f + feat_env + feat_wpt + feat_stft + feat_hht + feat_bearing
            all_features.append(feature_vec)

        labels.extend([full_label] * len(segments))  # 用 list 累加

    # full_label = parse_label_full(file_path, ch)  # 生成复合标签
    # print(full_label)
    # labels = [full_label] * len(all_features)

    return np.array(all_features, dtype=object), np.array(labels)

# ================= 遍历数据集 =================
def extract_dataset_features(root_dir, fs=12000, overlap=0.5):
    all_features, all_labels = [], []
    for root, dirs, files in os.walk(root_dir):
        for file in files:
            if file.endswith(".mat"):
                file_path = os.path.join(root, file)
                feats, labs = extract_features_from_file(file_path, fs, overlap)
                all_features.append(feats)
                all_labels.append(labs)
                print(f"Processed {file_path}, {feats.shape[0]} samples.")
    X = np.vstack(all_features)
    y = np.hstack(all_labels)
    return X, y

# ================= 保存 CSV =================
def save_features_to_csv(X, y, save_path="bearing_features_full.csv"):
    # 时域特征
    time_cols = ["mean","std","var","RMS","peak","crest_factor","impulse_factor","margin_factor","kurtosis","skew"]
    # 频域特征
    freq_cols = ["freq_centroid","peak_freq","bandwidth","spectral_kurtosis"]
    # 包络谱
    env_cols = ["env_energy","env_peak_freq","env_spectral_kurt"]
    # 小波包
    wpt_cols = []
    for i in range(8):
        wpt_cols += [f"wpt_energy_{i}", f"wpt_ratio_{i}"]
    wpt_cols += ["wpt_entropy"]
    # STFT
    stft_cols = ["stft_mean","stft_std","stft_max"]
    # HHT
    hht_cols = ["hht_amp_mean","hht_amp_std","hht_freq_mean","hht_freq_std"]
    # 轴承特征
    bearing_cols = ["fr","BPFO","BPFI","BSF","FTF"]

    columns = ["Channel"] + time_cols + freq_cols + env_cols + wpt_cols + stft_cols + hht_cols + bearing_cols

    df = pd.DataFrame(X, columns=columns)
    df["Label"] = y
    df.to_csv(save_path, index=False)
    print(f"Saved features to {save_path}")

# ================= 主函数 =================
if __name__ == "__main__":
    root_dir = "/home/xxxten/PycharmProjects/jianmo/data/E题数据集/源域数据集/48kHz_DE_data"
    X, y = extract_dataset_features(root_dir, fs=48000, overlap=0.5)
    # X, y = extract_dataset_features(root_dir, fs=12000, overlap=0.0)
    print("Final dataset shape:", X.shape, y.shape)
    save_features_to_csv(X, y, save_path="../data/feature_1/bearing_features_full_48kHz_DE_data")
\end{Python}
\vspace{2ex}

% \clearpage
